\documentclass[11pt]{article}
\usepackage[T1]{fontenc}
\usepackage{amsmath,amssymb}
\usepackage{fourier}
\usepackage{booktabs}
\usepackage{enumitem}
\usepackage[margin=1in]{geometry}
\usepackage[hidelinks]{hyperref}
\usepackage{fancyvrb}

\newcommand{\code}[1]{\texttt{#1}}
\setlength{\emergencystretch}{3em}
\fvset{fontsize=\small,xleftmargin=1.5em}

\title{Installing and using LLR with \texttt{mepn-python}}
\author{}  % add author
\date{21 November 2026, 20:52 MST (UTC$-$07:00)}

\begin{document}
\maketitle

\begin{abstract}
LLR, by Jean Penn\'e, tests numbers of the form $k\cdot b^n\pm c$ and their quotients for primality, using fast arithmetic for these special forms. For long searches it is much faster than GMP, and for forms such as $k\cdot b^n\pm1$ it proves primality instead of testing for probable primality. This report explains how to install it under Linux or WSL, how \code{search} finds and calls it, how to check that it is in use, and when it is worth having.
\end{abstract}

\section{Getting LLR}

LLR is distributed as ready-made executables from Jean Penn\'e's site,
\begin{center}
\url{http://jpenne.free.fr/index2.html}.
\end{center}
For Linux, including WSL under Windows, download the 64-bit Linux build. It is usually a zip archive containing a statically linked executable named \code{sllr64}; being static, it needs no further libraries. The current version and the exact file names were not checked for this report, so use whatever the page offers for 64-bit Linux.

\section{Installing it}

\begin{enumerate}[label=\arabic*.]
\item Download the zip in the Windows browser. WSL sees the Windows Downloads folder under \code{/mnt/c}:
\begin{Verbatim}
cd ~ && mkdir -p llr && cd llr
cp "/mnt/c/Users/<Windows name>/Downloads/<llr zip file>" .
unzip "<llr zip file>"           # sudo apt install unzip, if needed
\end{Verbatim}

\item Copy the executable to the place where \code{search} looks first, and make it executable. Unzipping often drops that permission.
\begin{Verbatim}
cp sllr64 ~/mepn-python/bin/linux-x64/sllr64
chmod +x ~/mepn-python/bin/linux-x64/sllr64
\end{Verbatim}
If the executable in the archive has another name, either keep it, since \code{search} also looks for \code{llr64} and \code{llr}, or rename it to \code{sllr64}.

\item Check that it runs:
\begin{Verbatim}
~/mepn-python/bin/linux-x64/sllr64 -v
\end{Verbatim}
\end{enumerate}

\paragraph{Windows.} A Windows build works too, but from Windows rather than from WSL: place \code{cllr64.exe} in \code{bin/windows-x64/}, together with its DLLs. Under WSL, use the Linux build.

\section{How \texttt{search} finds and uses it}

\paragraph{Finding it.} Unless told otherwise, \code{search} looks for an executable named \code{sllr64}, \code{cllr64}, \code{llr64} or \code{llr}, first in \code{bin/<platform>/} of the repository, then in the current folder, then on the \code{PATH}. It can also be named directly:
\begin{Verbatim}
python -m bright search 300 10000 --llr /path/to/sllr64
export BRIGHT_LLR=/path/to/sllr64          # the same, for every run
python -m bright search 300 10000 --no-llr  # force the built-in GMP test
\end{Verbatim}
The first line printed by \code{search} names the test in use, for example \code{primality test: /root/mepn-python/bin/linux-x64/sllr64} or \code{primality test: built-in}. If an LLR file exists but is not executable, \code{search} says so and suggests \code{chmod +x}.

\paragraph{Calling it.} For each candidate, \code{search} writes a one-line input file in LLR's ABC format and runs LLR on it with \code{-d}:
\begin{Verbatim}
ABC ($a*$b^$c$d)/$e
4105 17 100000 -9 16
\end{Verbatim}
This line stands for $(4105\cdot17^{100000}-9)/16$, a member of the base-17 family \code{F1\{9\}}. In general, a member of a simple family $x\,y^*z$ is $(A\cdot b^{c}+C)/q$, where the exponent $c$ is the number of repeated digits plus the length of $z$ (\code{bright-algorithm.tex}, Section~11). An optional fifth argument to \code{search} is passed on to LLR as \code{-a<arg>}, as in the C version.

\paragraph{Reading the result.} \code{search} treats ``not prime'' as composite, and both ``PRP'' and ``is prime'' as a prime. LLR reports the latter when it can prove primality, as it can for forms like $k\cdot b^n\pm1$. The C program recognised only ``PRP'', and would have retried a proven prime forever. If LLR's output contains neither, \code{search} stops with the message \code{LLR gave no result}, rather than guess.

\paragraph{Its own files.} LLR writes files such as \code{lresults.txt} and \code{llr.ini} in the folder it is run from. They are harmless; this is one more reason to use a separate folder for each run.

\section{Checking the installation}

Run a small, known case and compare the result with the stored set:
\begin{Verbatim}
mkdir -p ~/runs/llrtest && cd ~/runs/llrtest
~/mepn-python/run-class.sh     # base 16, class 3 mod 4,
                               # simple to 100, search 100 to 2000
grep "primality test" 16.3mod4/run.log
diff <(sort 16.3mod4/data/minimal.16.3mod4.txt) \
     <(sort ~/mepn-python/data-computed/minimal.16.3mod4.txt) && echo same
\end{Verbatim}
The log should name LLR, and the 359 primes should match. The last family of this class, \code{F8*F}, gives the prime $\mathrm{F}\,8^{1517}\mathrm{F}$, which LLR has to find.

\section{When it is worth it}

The built-in test calls GMP's \code{mpz\_probab\_prime\_p}. On random composites without small factors it took

\begin{center}
\begin{tabular}{@{}lrrrrrr@{}}
\toprule
digits & 1{,}000 & 2{,}000 & 4{,}000 & 8{,}000 & 16{,}000 & 32{,}000\\
seconds & 0.008 & 0.05 & 0.28 & 1.6 & 8.7 & 45\\
\bottomrule
\end{tabular}
\end{center}

which grows like $d^{2.5}$ in the number of digits (\code{complexity.tex}). LLR uses arithmetic specialised to numbers of this form, and its advantage grows with size. LLR was not available where these measurements were made, so no comparison is given here.
\begin{itemize}
\item For searches up to a few thousand digits, the GMP test is adequate, and most of the class sets in \code{data-computed} were completed with it.
\item For the open families searched to $20{,}000$ digits and beyond, such as base 10, $1 \bmod 3$ (\code{8*9}) or base 12, $1 \bmod 8$ (\code{10*1}), and for the unsolved families of bases 17--29, LLR is the right tool.
\item For forms $k\cdot b^n\pm1$, LLR also turns a probable prime into a proven one.
\end{itemize}
Its output is the same either way: a family's first prime, which is then added to the kernel.

\end{document}
